\documentclass[10pt,a4paper]{amsart}
\usepackage[margin=19mm]{geometry}
\usepackage{amsmath,amssymb,mathtools}
\usepackage[scr=rsfs,cal=euler]{mathalfa}
\usepackage{xfrac}
\usepackage[unicode,hidelinks,bookmarks=false]{hyperref}
\newcommand{\ie}{i.e.\ }
\newcommand{\cf}{cf.\ }
\newcommand{\resp}{respectively\ }
\newcommand{\Subsection}{\subsection}
\providecommand{\nobreakdash}{\nobreak\hbox{-}}
\setlength{\emergencystretch}{2em}
\multlinegap0pt
\usepackage{bm}
\usepackage{csquotes}
\usepackage{amssymb}
\newtheorem{lemma}{Lemma}
\title{Limiting Distribution and Rate of Convergence for GL(3) Fourier Coefficients}
\author{Zongqi Yu}
\date{Source: arXiv:2605.20707v1 (CC BY 4.0)}
\begin{document}
\maketitle

\noindent\emph{Source and licence.} Limiting Distribution and Rate of Convergence for GL(3) Fourier Coefficients. Version arXiv:2605.20707v1 (CC BY 4.0), distributed by arXiv under the Creative Commons Attribution 4.0 International licence, http://creativecommons.org/licenses/by/4.0/, as recorded in the arXiv licence metadata for this exact version and shown on the abstract page https://arxiv.org/abs/2605.20707v1. Full licence notice: \texttt{licenses/arxiv-2605.20707-CC-BY-4.0.txt}.\par\smallskip

\noindent\textit{Editorial scope.} Counted unit: the article's lemma le6.5 (source0.tex lines 905--914). The article's complete original proof of it is delivered with it (source0.tex lines 915--1059, one \texttt{proof} environment of 145 lines). No mathematics is written, repaired or re-derived: the statement and the proof are the article's own lines, in the article's own notation, and no retained line is substituted or re-flowed.  Retained uncounted as the source-local context the proof needs: lines 726--734; lines 807--815; lines 863--872.  Nothing was omitted from any retained span, so the package carries no omission record.  No retained line contains a citation, so the wrapper carries no bibliography and no bibliography entry.  No retained line contains a figure environment, an image-inclusion command or drawing code, so no image is delivered and the package declares no asset dependency.  Editorial formatting only, in addition: an AMS \texttt{amsart} preamble that restates the article's own declarations and macros used by the retained lines, namely the environments \texttt{lemma} on separate counters, together with the standard packages those lines need.  The retained environments print in this document as Lemma 1 in order of appearance, whereas the article prints the same items with the numbering of its own full document.  The complete native source of the article as distributed in the arXiv e-print for this exact version is bundled unchanged as \texttt{sources/201012/source0.tex}.
\begin{lemma}\label{le6.1}
    Let $f$ be a fixed self-dual Hecke--Maass cusp form for $\operatorname{GL}(3)$. Let $\beta_0$ be a real constant. Let $\left\{a_n\right\}_n$ be a sequence of real numbers such that $\left|a_n\right| \leqslant 1$. There exist positive constants $c_{f,1}, c_{f,2}$ such that for any positive integers $1 \leqslant N<M, L \geqslant 2$ and $k \geqslant 1$ we have
    \begin{equation}\label{eq6.1}
    \begin{aligned}
        & \mathbb{E}\left(\left|\sum_{N \leqslant n<M} \frac{\varepsilon(n)}{n^{2 / 3}} \sum_{q \leqslant L} a_{n q^3} \frac{A_f\left(1,n q^3\right)}{q^{2}} \cos \left(6 \pi q \mathbb{X}_n\right)\right|^k\right) \\
        & \leqslant \min \left\{\left(c_{f,1} k^{2 / 3}\left(\log 2k\right)^{2/3}\right)^k,\left(c_{f,2} kN^{-1/3}\right)^{k / 2}\right\}.
    \end{aligned}
    \end{equation}
\end{lemma} 

\begin{lemma}\label{le6.3}
     Let $a_m$ be real numbers such that $a_m \ll 1$. Let $\alpha_0 \neq 0$ and $\beta_0$ be fixed real numbers. Let $T$ be large, $0 \leqslant h \leqslant(\log \log T) / 4$ be an integer and $M \leqslant T^{1 /\left(3^h+9 h\right)} / h^3$ be a real number. Write $m=n r^3$ where $n$ is cube-free and let $\left\{\mathbb{X}_n\right\}_{n \text { cube-free  }}$be a sequence of independent random variables uniformly distributed on $[0,1]$. Then we have
     \[
     \begin{aligned}
     &\frac{1}{T} \int_T^{2 T}\left(\sum_{m \leqslant M} a_m \cos \left(6 \pi \alpha_0 \sqrt[3]{m t}\right)\right)^h d t\\
     &=  \mathbb{E}\left(\left(\sum_{n r^2 \leqslant M} a_{n r^2} \cos \left(6 \pi r \mathbb{X}\right)\right)^h\right) +O\left(T^{- 2/ 9}\right).
     \end{aligned}
     \]
\end{lemma}

\begin{lemma}\label{le6.4}
    Let $T$ be large and let $f$ be a fixed self-dual Hecke--Maass cusp form for $\operatorname{GL}(3)$. For a positive integer $N$, we define
\[
F_{f,N}(t):=\frac{1}{\pi \sqrt{3}} \sum_{n \leqslant N} \frac{\varepsilon(n)}{n^{2/3}} \sum_{r \leqslant N} \frac{A_f\left(1,n r^3\right)}{r^{ 2}} \cos \left(6 \pi r \sqrt[3]{n t}\right).
\]
If $N \leqslant T^{1 /162}$ then for any fixed $\varepsilon>0$ we have
\[
\frac{1}{T} \int_T^{2 T}\left|F_f(t)-F_{f,N}(t)\right| d t \ll \frac{1}{N^{1 / 8}}.
\]
\end{lemma}

\begin{lemma}\label{le6.5}
    Let $f$ be a fixed self-dual Hecke--Maass cusp form for $\operatorname{GL}(3)$. There exists  a set $\mathcal{A} \in[T, 2 T]$ verifying $\operatorname{meas}([T, 2 T] \backslash \mathcal{A}) \ll T(\log T)^{-20}$ and a positive constant $c_{f,0}$, such that for all complex numbers $\lambda$ with
    \[
    |\lambda| \leqslant c_{f,0}\frac{\left(\log\log T\right)^{1/3}}{\left(\log\log\log T\right)^{5/3}},
    \] 
    we have
    \[
    \frac{1}{T} \int_{\mathcal{A}} \exp (\lambda F_f(t)) d t=\mathbb{E}\left(e^{\lambda F_{f,\mathbb{X}}}\right)+O\left(\exp \left(-\frac{\log \log  T}{\log\log \log  T}\right)\right).
    \]
\end{lemma}

\begin{proof}
Let $N: = \exp(\sqrt{\log T})$. Denote by $\mathcal{A}_1$ the set of all $t \in [T, 2T]$ for which
\begin{equation}\label{eq6.4}
\bigl| F_f(t) - F_{f,N}(t) \bigr| \le N^{-1/20}.
\end{equation}
From Lemma \ref{le6.4} we obtain
\begin{equation}\label{eq6.5}
    \operatorname{meas}\left([T, 2 T] \backslash \mathcal{A}_1\right) \leq N^{1 / 20} \int_T^{2 T}\left|F_f(t)-F_{f,N}(t)\right| d t \ll \frac{T}{N^{1 / 20}}.
\end{equation}
Let $K := \bigl\lfloor (\log\log T)/8 \bigr\rfloor$ and define $\mathcal{A}_2$ as the set of $t \in [T, 2T]$ satisfying
\begin{equation}\label{eq6.6}
\bigl| F_{f,N}(t) \bigr| \le V,
\end{equation}
where $V = c_{f,3} K^{2/3}\left(\log 2K\right)^{2/3}$ for a suitably large constant $c_{f,3}$. We also set
\[
F_{f,N, \mathbb{X}}:=\frac{1}{\pi \sqrt{3}} \sum_{n \leqslant N} \frac{\varepsilon(n)}{n^{2 / 3}} \sum_{r \leqslant N} \frac{A_f\left(1,n r^3\right)}{r^{ 2}} \cos \left(6\pi r \mathbb{X}_n\right),
\]
where $\left\{\mathbb{X}_n\right\}_{n \text { cube-free }}$ is a sequence of independent random variables, uniformly distributed on $[0,1]$. Then it follows from Lemma \ref{le6.3} that for any integer $0 \leqslant k \leqslant 2 K$ we have
\begin{equation}\label{eq6.7}
    \frac{1}{T} \int_T^{2 T} F_{f,N}(t)^k d t=\mathbb{E}\left(F_{f,N, \mathbb{X}}^k\right)+O\left(T^{-2 / 9}\right).
\end{equation}
Combining this and  Lemma \ref{le6.1}, we have
\begin{equation}\label{eq6.8}
    \begin{aligned}
\operatorname{meas}\left([T, 2 T] \backslash \mathcal{A}_2\right) & \leqslant V^{-2 K} \int_T^{2 T}\left|F_{f,N}(t)\right|^{2 K} d t \\
& \ll V^{-2 K} T \cdot \mathbb{E}\left(\left|F_{f,N, \mathrm{X}}\right|^{2 K}\right)+V^{-2 K} T^{7 / 9} \\
& \ll T\left(\frac{2 c_{f,1} K^{2 /3}\left(\log 2K\right)^{2/3}}{V}\right)^{2 K} \ll \frac{T}{(\log T)^{20}}.
\end{aligned}
\end{equation}
with a sufficiently large $c_{f,3}$. Define $\mathcal{A} := \mathcal{A}_1 \cap \mathcal{A}_2$ and  $\mathcal{A}^c = [T, 2T] \setminus \mathcal{A}$. Combining \eqref{eq6.5} and \eqref{eq6.8}, we obtain
\begin{equation}\label{eq6.9}
    \operatorname{meas}\left(\mathcal{A}^c\right) \ll \frac{T}{(\log T)^{20}} .
\end{equation}
Now by \eqref{eq6.4} and 
\[
\frac{1}{T} \int_{\mathcal{A}} \exp \left(\operatorname{Re}(\lambda) F_{f,N}(t)\right) d t \leqslant \exp (\operatorname{Re}(\lambda) V) \ll N^{1 / 50},
\]
we get
\begin{equation}\label{eq6.10}
    \begin{aligned}
\frac{1}{T} \int_{\mathcal{A}} \exp (\lambda F_f(t)) d t & =\frac{1}{T} \int_{\mathcal{A}} \exp \left(\lambda F_{f,N}(t)+O\left(|\lambda| N^{-1 / 20}\right)\right) d t \\
& =\frac{1}{T} \int_{\mathcal{A}} \exp \left(\lambda F_{f,N}(t)\right) d t+O\left(N^{-1 / 30}\right).
\end{aligned}
\end{equation}
We now estimate the main term on the right‑hand side of above. Set $L := \bigl\lfloor \frac{\log \log T}{\log \log \log T}  \bigr\rfloor$. By Stirling’s formula and the assumption on $\mathcal{A}$, we have
\begin{equation}\label{eq6.11}
    \frac{1}{T} \int_{\mathcal{A}} \exp \left(\lambda F_{f,N}(t)\right) d t=\sum_{k=0}^{2 L} \frac{\lambda^k}{k!} \frac{1}{T} \int_{\mathcal{A}} F_{f,N}(t)^k d t+E_1.
\end{equation}
where
\[
E_1 \ll \sum_{k>2 L} \frac{(|\lambda| V)^k}{k!} \ll \sum_{k>2 L}\left(\frac{3|\lambda| V}{k}\right)^k \ll \sum_{k>2 L}\left(\frac{3|\lambda| V}{2 L}\right)^k \ll e^{-L},
\]
with $c_{f,0}$ is sufficiently small. By the Cauchy--Schwarz inequality, Lemma \ref{le6.1}, \eqref{eq6.7} and \eqref{eq6.9} gives, for every $k \leqslant 2L$,
\[
\begin{aligned}
\left|\int_{\mathcal{A}^c} F_{f,N}(t)^k d t\right| & \leqslant \operatorname{meas}\left(\mathcal{A}^c\right)^{1 / 2}\left(\int_T^{2 T}\left|F_{f,N}(t)\right|^{2 k} d t\right)^{1 / 2} \\
& \ll \frac{\sqrt{T}}{(\log T)^{10}}\left(T \cdot \mathbb{E}\left(\left|F_{f,N, \mathrm{X}}\right|^{2 k}\right)+T^{7 / 9}\right)^{1 / 2} \\
& \ll \frac{T}{(\log T)^{10}}\left(2 c_{f,1} k^{2 / 3}\left(\log 2k\right)^{2/3}\right)^k \ll \frac{T}{(\log T)^9} .
\end{aligned}
\]
Inserting this estimate in \eqref{eq6.11} and using  \eqref{eq6.7} we get
\begin{equation}\label{eq6.12}
    \frac{1}{T} \int_{\mathcal{A}} \exp \left(\lambda F_{f,N}(t)\right) d t=\sum_{k=0}^{2 L} \frac{\lambda^k}{k!} \mathbb{E}\left(F_{f,N, \mathbb{X}}^k\right)+E_2,
\end{equation}
where
\[
E_2 \ll e^{-L}+\frac{1}{(\log T)^9} \sum_{k=0}^{2 L} \frac{|\lambda|^k}{k!} \ll e^{-L}+\frac{e^{|\lambda|}}{(\log T)^9} \ll e^{-L} .
\]
Combining above with \eqref{eq6.10} we derive
\begin{equation}\label{eq6.13}
    \frac{1}{T} \int_{\mathcal{A}} \exp (\lambda F_f(t)) d t=\sum_{k=0}^{2 L} \frac{\lambda^k}{k!} \mathbb{E}\left(F_{f,N, \mathbb{X}}^k\right)+O\left(e^{-L}\right).
\end{equation}
It remains to prove that the main term on the right‑hand side of \eqref{eq6.13} is asymptotically close to $\mathbb{E}\bigl(e^{\lambda F_{\mathbb{X}}}\bigr)$. For this purpose, define the event $\mathcal{B}$ by
\[
\left|F_{f,\mathrm{X}}-F_{f,N, \mathrm{X}}\right| \leqslant N^{-1 / 20} .
\]
Then, it follows from Lemma \ref{le6.1} that
\begin{equation}\label{eq6.14}
    \begin{aligned}
\mathbb{P}\left(\mathcal{B}^c\right) & \leqslant N^{1 / 10} \mathbb{E}\left(\left|F_{f,\mathbb{X}}-F_{f,N, \mathbb{X}}\right|^2\right) \\
& \ll N^{1 / 10} \mathbb{E}\left(\left|\sum_{n>N} \frac{\varepsilon(n)}{n^{2/3}} \sum_{r=1}^{\infty} \frac{A_f\left(1,n r^3\right)}{r^{2}} \cos \left(6 \pi r \mathbb{X}_n\right)\right|^2\right)+N^{-1 /5} \\
& \ll N^{-1 / 5},
\end{aligned}
\end{equation}
since, noting that
\[
\sum_{n\leqslant N} \sum_{r> N}\frac{\varepsilon(n)^2A_f(1,nr^3)^2}{\left(nr^3\right)^{4/3}}\leqslant \sum_{n>N}\frac{A_f(1,n)^2}{n^{4/3}}, 
\]
we have
\[
 \mathbb{E}\left(\left|\sum_{n\leqslant N} \frac{\varepsilon(n)}{n^{2/3}} \sum_{r>N}^{\infty} \frac{A_f\left(1,n r^3\right)}{r^{2}} \cos \left(6 \pi r \mathbb{X}_n\right)\right|^2\right)
 \ll N^{-1 /3} .
\]

Let $\mathbf{1}_{\mathcal{B}}$ denote the indicator function of an event $\mathcal{B}$, and let $k \leqslant 2L$ be a positive integer. We then note that
\[
\begin{aligned}
\mathbb{E}\left(F_{f,\mathbb{X}}^k\right) & =\mathbb{E}\left(\mathbf{1}_{\mathcal{B}} \cdot F_{f,\mathbb{X}}^k\right)+\mathbb{E}\left(\mathbf{1}_{\mathcal{B}^c} \cdot F_{f,\mathbb{X}}^k\right) \\
& =\mathbb{E}\left(\mathbf{1}_{\mathcal{B}} \cdot\left(F_{f,N, \mathbb{X}}+O\left(N^{-1 / 20}\right)^k\right)+\mathbb{E}\left(\mathbf{1}_{\mathcal{B}^c} \cdot F_{f,\mathbb{X}}^k\right)\right. \\
& =\mathbb{E}\left(\left(F_{f,N, \mathbb{X}}+O\left(N^{-1 / 20}\right)\right)^k\right)\\
&+O\left(\mid \mathbb{E}\left(\mathbf{1}_{\mathcal{B}^c} \cdot\left(F_{f,N, \mathbb{X}}+O\left(N^{-1 / 20}\right)^k\right)\left|+\left|\mathbb{E}\left(\mathbf{1}_{\mathcal{B}^c} \cdot F_{f,\mathbb{X}}^k\right)\right|\right.\right.\right) .
\end{aligned}
\]
 Applying the Binomial Theorem together with Lemma \ref{le6.1}, we obtain
\[
\mathbb{E}\left(F_{f,N, \mathbb{X}}^k\right)+E_3,
\]
where
\[
\begin{aligned}
    E_3 &\ll \sum_{j=1}^k\binom{k}{j}\left(c_{f,4} N\right)^{-j / 20} \mathbb{E}\left(\left|F_{f,N, \mathrm{X}}\right|^{k-j}\right)\\
    &\ll N^{-1 / 20}\left(c_{f,5} k^{2 / 3}\left(\log 2k\right)^{2/3}\right)^k \ll N^{-1 / 40},
\end{aligned}
\]
for some positive constant $c_{f,4}$ and $c_{f,5}$. Next, by the Cauchy--Schwarz inequality, Lemma \ref{le6.1} and  \eqref{eq6.14} we deduce that
\[
\begin{aligned}
& \mid \mathbb{E}\left(\boldsymbol{1}_{\mathcal {B}^{c}}\cdot (F_{f, N,\mathbb{X}}+O(N^{-1/20} )^{k} ) \left|+\left|\mathbb{E}\left(\mathbf{1}_{\mathcal{B}^c} \cdot F_{f,\mathbb{X}}^k\right)\right|\right.\right. \\
& \leqslant \mathbb{P}\left(\mathcal{B}^c\right)^{1 / 2}\left(\mathbb{E}\left(\left|F_{f,N, \mathbb{X}}+O\left(N^{-1/20}\right)\right|^{2 k}\right)^{1 / 2}+\left|\mathbb{E}\left(\left|F_{f,\mathbb{X}}\right|^{2 k}\right)\right|^{1 / 2}\right) \\
& \ll N^{-1 / 10}\left(2c_{f,1} k^{2 / 3}\left(\log 2k\right)^{2/3}\right)^k \ll N^{-1 / 20} .
\end{aligned}
\]
Collecting the above estimates, we obtain
\[
\mathbb{E}\left(F_{f,N, \mathbb{X}}^k\right)=\mathbb{E}\left(F_{f,\mathbb{X}}^k\right)+O\left(N^{-1 / 40}\right) .
\]
Substituting the above into \eqref{eq6.13} yields
\[
\begin{aligned}
\frac{1}{T} \int_{\mathcal{A}} \exp (\lambda F_f(t)) d t & =\sum_{k=0}^{2 L} \frac{\lambda^k}{k!} \mathbb{E}\left(F_{f,\mathbb{X}}^k\right)+O\left(N^{-1 / 40} e^{|\lambda|}+e^{-L}\right) \\
& =\mathbb{E}\left(e^{\lambda F_{f,\mathbb{X}}}\right)+E_4,
\end{aligned}
\]
where
\[
E_4 \ll e^{-L}+\sum_{k>2 L} \frac{|\lambda|^k}{k!} \mathbb{E}\left(\left|F_{f,\mathbb{X}}\right|^k\right) \ll e^{-L}+\sum_{k>2 L}\left(\frac{3 c_{f,1}|\lambda|\left(\log 2k\right)^{2/3}}{k^{1 /3}}\right)^k,
\]
by Stirling's formula and Lemma \ref{le6.1}. Finally, our assumption on $\lambda$ ensures that for the tail of the series,
\[
\sum_{k > 2L} \left( \frac{3 c_{f,1} |\lambda|\left(\log 2k\right)^{2/3}}{k^{1 /3}} \right)^k
\ll \sum_{k > 2L} \left( \frac{3 c_{f,1} |\lambda|\left(\log 4L\right)^{2/3}}{L^{1 /3}} \right)^k
\ll e^{-L}.
\]
 This completes the proof.
\end{proof}

\end{document}
